% Related work section for tech report

Several lines of prior work are directly relevant. Twelve papers were
reviewed in full in the project literature assessment, and six additional
works received abstract-level screening. The reviewed set is targeted,
not an exhaustive search of the field.

\subsection{Finite-Word-Length FIR Design}

Kodek and Steiglitz~\cite{kodek1980filter} analyzed length-wordlength
tradeoffs for FIR filters, proving a nonzero lower bound on weighted
Chebyshev error for fixed coefficient wordlength. Their complexity proxy
$Nb$ motivates joint optimization; their analysis assumes a common
wordlength, not per-coefficient quotas.

Lim, Parker, and Constantinides~\cite{lim1982finite} minimized FIR
designs over discrete coefficient-value spaces using integer linear and
quadratic programming, comparing discrete-space designs with rounded
infinite-wordlength designs. Lim~\cite{lim1990nprm} extended this to
normalized peak-ripple minimization, using MILP methods with passband
gain as a design variable.

\subsection{SPT-Based and Hardware-Aware Coefficient Allocation}

Several methods allocate signed-power-of-two (SPT) terms across
coefficients. Llorens, Hadjicostis, and Ni~\cite{llorens2005integercost}
formulated FIR coefficient quantization under a scalar total integer-cost
constraint using a constrained-shortest-path DP. Their alpha=2 case
relates coefficient-vector error to frequency-response MSE via Parseval's
theorem. For exact three-width quotas, a radix encoding
($r_{16}=B+1$, $r_8=2$, $r_4=1$, $B=N+1$) reduces the quota DP to a
scalar exact-cost DP, so the present quota DP is a specialized
realization of Llorens et al.'s framework, not a new optimization class.
They also describe a greedy algorithm with a conditional optimality
guarantee: the guarantee depends on the marginal error reduction as the
integer cost increases. The sparse costs used in the radix encoding do not
automatically satisfy that condition: intermediate cost levels represent no
additional width choice, while the final jump to 16 bits can have a positive
error reduction. The greedy baseline in this report instead selects the
largest immediate squared-error reduction among quota-feasible 4-to-8 and
8-to-16 upgrades. It is a comparator only; no optimality guarantee is
claimed, and it is not identified with Llorens et al.'s conditionally
optimal greedy algorithm.

Chen and Willson~\cite{chen1999trellis} used a trellis (Viterbi-like)
search for SPT term allocation against frequency-response criteria; their
survivor pruning does not guarantee global optimality. Li, Lim, Lian,
and Song~\cite{li2002polynomial} allocated a global SPT-term budget with
variable per-coefficient counts using greedy sequential allocation;
the authors describe the result as suboptimal.

Aktan, Yurdakul, and D\"{u}ndar~\cite{aktan2008algorithm} used
branch-and-bound (FIRGAM) with LP for coefficient bounds to reduce SPT
count while meeting frequency specifications; worst-case runtime is
exponential.

Kuo, Lin, and Liu~\cite{kuo2011complexity} integrated successive SPT
allocation with common-subexpression elimination under an exact global
post-CSE addition budget, reporting 0.35-$\mu$m CMOS synthesis results.

\subsection{Subexpression-Space and Structure Optimization}

Yu and Lim~\cite{yu2007subexpression} optimized FIR coefficients
directly in a discrete subexpression space using branch-and-bound MILP,
minimizing frequency-response ripple subject to an adder budget. This is
a direct coefficient-space search, not a two-stage quantize-then-CSE
pipeline.

Blad and Gustafsson~\cite{blad2010integer} used ILP to optimize
partial-product reduction trees and pipeline registers for fixed
coefficient representations, targeting 90-nm CMOS. Their decision
variables are hardware structural elements, not coefficient widths.

\subsection{Mixed-Precision FIR Allocation}

Fang, Chen, Chen, Wang, and You~\cite{fang2025mixedprecision} proposed a
PSO-based bit-allocation framework applied to FIR filter design with a
frequency-domain minimax objective and aggregate bit-consumption
constraint. Their constraint is an inequality on total bits ($\sum b_n
\le N \bar{b}$), not exact per-width counts. Their allocation assumes
symmetric widths (linear-phase structure). This is the closest prior
mixed-precision FIR allocation paper; it differs from the present work
in constraint type, objective, and solution method.

Sarma, Sai, and Chatterjee~\cite{sarma2026coefficient} used a
coefficient-magnitude threshold to select exact or LUT-approximate
multipliers under an output-error budget on a Zynq UltraScale+ radar DDC
implementation. Their quantization model (approximate arithmetic)
differs from fixed-point coefficient widths.

\subsection{Neural-Network Mixed-Precision}

Yang and Jin's FracBits~\cite{yang2021fracbits} assigns fractional
bit-widths in neural-network quantization. It is included as
methodological context; it is not a direct FIR baseline.
